\documentclass[11pt,a4paper]{article}

\usepackage[utf8]{inputenc}
\usepackage[T1]{fontenc}
\usepackage{lmodern}
\usepackage[margin=0.76in]{geometry}
\usepackage[table]{xcolor}
\usepackage{array}
\usepackage{tabularx}
\usepackage{booktabs}
\usepackage{enumitem}
\usepackage{titlesec}
\usepackage{fancyhdr}
\usepackage{lastpage}
\usepackage{hyperref}

\definecolor{TextInk}{HTML}{233142}
\definecolor{HeadingBlue}{HTML}{1D3557}
\definecolor{AccentTeal}{HTML}{2A7F92}
\definecolor{AccentGold}{HTML}{C58B4E}
\definecolor{TableStripe}{HTML}{F3F8FA}
\definecolor{SoftBlue}{HTML}{EEF6F8}
\definecolor{SoftGreen}{HTML}{EEF7F0}
\definecolor{SoftRed}{HTML}{FBF0EF}

\hypersetup{
  colorlinks=true,
  linkcolor=HeadingBlue,
  urlcolor=AccentTeal,
  pdftitle={IB Geography SL May 4 Remediation Marking Guide},
  pdfauthor={IB Geography SL Preparation}
}

\color{TextInk}
\setlength{\parskip}{0.45em}
\setlength{\parindent}{0pt}
\setlength{\tabcolsep}{5pt}
\renewcommand{\arraystretch}{1.18}
\setlist[itemize]{leftmargin=1.32em,itemsep=3pt,topsep=4pt}

\newcolumntype{Y}{>{\raggedright\arraybackslash}X}
\newcolumntype{L}[1]{>{\raggedright\arraybackslash}p{#1}}

\titleformat{\section}
  {\normalfont\Large\sffamily\bfseries\color{HeadingBlue}}
  {}{0pt}{}
  [\vspace{0.08em}{\color{AccentGold}\titlerule[0.8pt]}]

\titleformat{\subsection}
  {\normalfont\large\sffamily\bfseries\color{AccentTeal}}
  {}{0pt}{}

\pagestyle{fancy}
\setlength{\headheight}{14pt}
\fancyhf{}
\fancyhead[L]{\sffamily\color{AccentTeal}May 4 Marking Guide}
\fancyhead[R]{\sffamily\color{HeadingBlue}IB Geography SL}
\fancyfoot[C]{\sffamily\color{HeadingBlue}\thepage\ of \pageref{LastPage}}
\renewcommand{\headrulewidth}{0.4pt}

\newcommand{\term}[1]{\textcolor{HeadingBlue}{\textbf{#1}}}
\newcommand{\exam}[1]{\textcolor{AccentTeal}{\textbf{#1}}}
\newcommand{\score}[1]{\textcolor{AccentTeal}{\textbf{#1}}}
\newcommand{\boxnote}[3]{%
\noindent\colorbox{#1}{%
  \begin{minipage}{0.965\textwidth}
  \sffamily
  \textbf{\textcolor{HeadingBlue}{#2}} #3
  \end{minipage}
}}

\begin{document}

\begin{center}
  {\sffamily\bfseries\color{HeadingBlue}\LARGE May 4 Remediation Marking Guide}\\[0.25em]
  {\sffamily\color{AccentTeal}\large Use After The Timed Redo Only}
\end{center}

\boxnote{SoftRed}{Strict rule.}{Mark only what was written during the timed
redo. Do not improve the answer before scoring. The purpose is to prove that
the May 2 or May 3 weakness has actually been fixed.}

\section{Task A Mark Guide: Source Evidence [2]}

Award \score{1 mark} for each valid piece of source evidence showing uneven
environmental quality. A full-mark answer should name a zone or category and
use an exact number or exact source condition.

Valid answers include:

\begin{itemize}
  \item The suburban zone has \term{34\%} green space, compared with only
  \term{3\%} in the peripheral informal settlement.
  \item The CBD has only \term{8\%} green space and faces congestion and
  limited open space.
  \item The peripheral informal settlement has weak service access and flood /
  pollution exposure.
  \item The inner-city regeneration zone has only \term{14\%} green space and
  faces redevelopment pressure.
\end{itemize}

\boxnote{SoftGreen}{Target.}{A corrected May 4 answer should score
\score{2/2}. If it scores \score{1/2}, the remaining problem is probably
missing data, missing category names, or no clear comparison.}

\section{Task B Mark Guide: Detroit 10-Mark Redo}

Use this best-fit level guide.

\rowcolors{2}{TableStripe}{white}
\begin{tabularx}{\textwidth}{L{0.14\textwidth}Y}
\toprule
\textbf{Marks} & \textbf{Descriptor} \\
\midrule
0 & No relevant answer. \\
1--2 & General statements about cities, decline, or gentrification with little
or no named evidence. \\
3--4 & Some relevant understanding of regeneration or gentrification, but
mainly descriptive and weakly supported. A named place may be mentioned but not
used well. \\
5--6 & Clear explanation with a named urban area and some balance. Evidence or
evaluation may be general. \\
7--8 & Good, balanced answer using Detroit or another named place effectively.
Includes clear benefits and limits, supported judgement, and some precise
case-study evidence. \\
9--10 & Detailed, well-structured evaluation with precise case-study evidence,
clear explanation of decline and regeneration, strong balance, and a direct
judgement about extent or success. \\
\bottomrule
\end{tabularx}
\rowcolors{2}{}{}

\subsection{Evidence To Reward}

Credit accurate use of:

\begin{itemize}
  \item Detroit, Michigan, United States as the named urban area;
  \item population decline from about \term{1.85 million} in 1950 to
  \term{639,111} in 2020;
  \item deindustrialization, automobile-sector restructuring, suburbanization,
  racial / income segregation, vacancy, and fiscal stress as causes of decline;
  \item selective reinvestment in downtown / Midtown, image improvement,
  service improvement, business activity, reuse of land, or job access as
  possible benefits;
  \item poverty, affordability pressure, displacement risk, exclusion, uneven
  benefits, or remaining service gaps as limits.
\end{itemize}

\subsection{Judgement To Reward}

Strong answers should judge degree. For example:

\begin{itemize}
  \item Regeneration partly solves urban decline by attracting investment and
  improving selected districts, but it does not fully solve city-wide poverty,
  affordability, or service inequality.
  \item Gentrification can improve the built environment and economic activity,
  but it is less socially sustainable if lower-income residents are excluded
  from the benefits.
  \item Detroit is best judged as uneven regeneration rather than a complete
  success story.
\end{itemize}

\section{Task C Mark Guide: Population Data [6]}

\subsection{C1: Two Alpha / Beta Contrasts [4]}

Award up to \score{4 marks}: up to \score{2 marks} for each clear contrast
using Source B evidence.

Valid contrasts include:

\begin{itemize}
  \item Country Alpha is more youthful: median age \term{18}, under-15 share
  \term{43\%}, and 65+ share \term{3\%}; Country Beta is ageing: median age
  \term{48}, under-15 share \term{12\%}, and 65+ share \term{29\%}.
  \item Alpha has much higher fertility (\term{4.6}) than Beta (\term{1.3}).
  \item Alpha is growing rapidly (\term{2.7\%}) while Beta is shrinking
  slightly (\term{-0.3\%}).
  \item Alpha has net out-migration (\term{-1.2\%}) while Beta has net
  in-migration (\term{0.4\%}).
\end{itemize}

\subsection{C2: Gamma Growth Near Replacement Fertility [2]}

Award up to \score{2 marks} for one explained reason. Good answers include:

\begin{itemize}
  \item Gamma has positive net migration of \term{0.8\%}, so immigration can
  add to total population even if fertility is near replacement.
  \item Population momentum may keep growth positive because previous larger
  cohorts are still moving through childbearing ages.
  \item If death rates are low or life expectancy is rising, total population
  may continue growing for a period.
\end{itemize}

\section{Final Decision}

\rowcolors{2}{TableStripe}{white}
\begin{tabularx}{\textwidth}{L{0.32\textwidth}Y}
\toprule
\textbf{Prompt} & \textbf{Fill in after marking} \\
\midrule
Corrected score & \rule{0.65\linewidth}{0.4pt} \\
Remaining error & \rule{0.65\linewidth}{0.4pt} \\
One sentence to carry into Paper 1 & \rule{0.65\linewidth}{0.4pt} \\
One Population data habit to carry into Paper 2 & \rule{0.65\linewidth}{0.4pt} \\
\bottomrule
\end{tabularx}
\rowcolors{2}{}{}

\end{document}
